\scanpage{36}
% archive sheet number: 35
\begingroup\makeatletter\tagsleft@true\makeatother
si \deleted{\unsure{Card|le}} cardinal de l'ens. $(18)$ est $4$, i.e. si les quatre parties de $X$
\[
\Sigma_1^+,\quad\Sigma_1^-,\quad\Sigma_2^+,\quad\Sigma_2^-
\]
sont distinctes. A ce propos notons que l'on sait déjà par \deleted{hypothèse $(\Pi_1,\Pi_2$ étant} \inserted{définition} des déc. binaires
\[
\Sigma_1^+\ne\Sigma_1^-,\qquad\Sigma_2^+\ne\Sigma_2^-
\tag*{($*$)}
\]
et qu'on a
\[
\Sigma_1^+\ne\Sigma_2^+
\tag*{($**$)}
\]
puisque $\Sigma_1^+$ et $\unsure{\Sigma^+|\Sigma_2^+}$ sont des parties disjointes \uline{non} \uline{vides}. Notons donc qu'on lit sur le diagramme $(14)$
\[
\left\{\begin{aligned}
\text{a)}\;(\Sigma_1^+\unsure{\ne|=}\Sigma_2^-)&\Longleftrightarrow E_1\unsure{\ne|=}\Sigma_{12}
   &(\Longrightarrow E_2=\varnothing)\\
\text{b)}\;(\Sigma_2^+=\Sigma_1^-)&\unsure{\Longleftrightarrow|\Longrightarrow} E_2=\Sigma_{12}
   &(\Longrightarrow E_1=\varnothing)
\end{aligned}\right.
\tag*{($***$)}
\]
et que les deux relations \inserted{a), b)} envisagées dans $(***)$ ne peuvent être satisfaites simultanément, puisque $\deleted{\Sigma_{12}}$ l'on aurait $\Sigma_{12}=E_1=E_2=\varnothing$ puisque $E_1\cap E_2=\varnothing$, ce qui contredit $(16)$ (c).

Enfin
\[
(\Sigma_1^-=\Sigma_2^-)\Longleftrightarrow
(\Sigma_{1\unsure{\deleted{2}|2}}^-=\unsure{\deleted{X|E}}\Sigma_2^-=X)
\quad\text{i.e. }\Sigma_{12}=X
\tag*{($****$)}
\]
Ce cas est mutuellement exclusif avec chacun des cas $(***)$, \deleted{\unsure{dans ces derniers}} \inserted{\unsure{Car|car} si on avait en $=$ temps} $\deleted{\unsure{\Sigma_1^-}=\Sigma_2^-=X}\ \deleted{\unsure{E}}\Sigma_{12}=X$ et $E_1=\Sigma_{12}$, on aurait $E_1=X$ donc $\Sigma_1^+=\Sigma_1^-=X$, \deleted{\unsure{de|et}}\unsure{ce} qui contredit $(3)$.

\uline{Proposition} Pour que l'on ait $\operatorname{Card}\{\Sigma_1^+,\Sigma_1^-,\Sigma_2^+,\Sigma_{\unsure{2|1}}^-\}=4$, il faut et il suffit que l'on ne soit pas l'un des trois cas suivants, qui s'excluent mutuellement
\[
\left\{\begin{aligned}
\text{\unsure{a})}\quad &E_1=\Sigma_{12}&&(\Longrightarrow E_2=\varnothing)\\
\text{a')}\quad &E_2=\Sigma_{12}&&(\Longrightarrow E_1=\varnothing)\\
\text{b)}\quad &\Sigma_{12}=X\quad\text{i.e. }\Sigma_1^-=\Sigma_2^-=X
   &&(\text{d'où }E_1=\Sigma_1^+,E_2=\Sigma_2^+)
\end{aligned}\right.
\tag*{(\deleted{\unsure{19|20}}20)}
\]
auquel cas le cardinal \unsure{du|des} $\unsure{\}\,4|\{\ \}|4}$ est égal à \unsure{trois}.

\endgroup
\clearpage
